\emph{Översikt}
Denna modul introducerar \emph{algoritmiskt tänkande}: att lösa ett problem
genom att beskriva lösningen som en följd av entydiga steg.  Vi utgår från
vardagliga algoritmer --- recept, tvättsortering, handdisk --- och urskiljer
ur dem en algoritms \emph{huvudbeståndsdelar}: variabler (det som kan
variera), funktioner (namngivna delalgoritmer), förgreningar (val mellan
alternativ) och upprepningar (att göra om steg).  Modulen avslutas med
\emph{stegvis förfining}: att utveckla en algoritm grovt först och sedan allt
mer detaljerat, tills stegen är så precisa att en dator kan följa dem.  På
vägen införs den notation med namngivna steg som kursens alla program
presenteras med.

\emph{Lärandemål}
Efter denna modul ska studenten kunna:
\begin{restatable}{lo}{CompThinkLOAlgorithm}\label{CompThinkLOAlgorithm}%
  Förklara vad en algoritm är --- en ändlig följd av entydiga steg som löser
  ett problem --- och ge exempel på algoritmer i vardagen.
\end{restatable}
\begin{restatable}{lo}{CompThinkLOComponents}\label{CompThinkLOComponents}%
  Urskilja och använda en algoritms huvudbeståndsdelar: variabler,
  funktioner, förgreningar och upprepningar.
\end{restatable}
\begin{restatable}{lo}{CompThinkLORefinement}\label{CompThinkLORefinement}%
  Utveckla en algoritm genom stegvis förfining av pseudokod, från grov
  överblick till körbara steg.
\end{restatable}
\begin{restatable}{lo}{CompThinkLONotation}\label{CompThinkLONotation}%
  Läsa notationen med namngivna steg (\LA{}namn\RA{} och
  \LA{}namn\RA{}\(\equiv\)) som kursens program presenteras med.
\end{restatable}

\emph{Förkunskaper}
Modulen förutsätter inga tidigare programmeringskunskaper.  Studenten bör kunna
följa och beskriva vardagliga stegvisa instruktioner, till exempel ett recept.

\emph{Läsning}
Kompletterande material finns i kursens videogenomgång \emph{Algoritmiskt
tänkande}.
